\section{Transport through finite trees of metrics}\label{sec:transport}

The relative metric estimate of Lemma~\ref{lem:relative-pin} measures the
entropy gained by one change of partition.  We now combine such changes
through a finite tree of matrix means.  The resulting derivative is an
average of entropy gains against positive, trace-one states.  The same
states control the energy of the filtered vector.  A second occurrence of
the probability that a term is split is essential in this energy estimate.
All trees, one-copy spaces, and partitions in this section are finite and
are fixed before the number of copies tends to infinity.
The matrix means carry perturbations from leaves to the root.  The trace
adjoints of their normalized derivatives carry a density matrix at the
root back to each leaf; this reverse passage will identify the states
used in both estimates.

\subsection{Positive matrix means and their derivatives}

For positive definite operators $A,B$ on a finite-dimensional Hilbert
space and $0\leq p\leq1$, define
\begin{equation}\label{transport:mean-def}
 A\#_p B=A^{1/2}(A^{-1/2}BA^{-1/2})^pA^{1/2}.
\end{equation}
A binary operation of this form is the weighted operator geometric mean
in the Kubo--Ando framework~\cite{KuboAndo1980}; the earlier geometric
mean of positive forms is due to Pusz and Woronowicz~\cite{PuszWoronowicz1975}.
We prove below the finite-tree properties and normalized derivative
identities needed for the entropy and energy estimates.
A weighted binary tree is a finite rooted binary tree with a parameter
$p\in[0,1]$ assigned to each internal vertex; its two outgoing edges
have weights $1-p$ and $p$, respectively.
Its terminal weights are the products of edge weights along root-to-leaf
paths; they sum to one.  Assign positive definite operators to the leaves
and evaluate each internal vertex by \eqref{transport:mean-def}.  This
defines the root operator.  An edge of weight zero is omitted, and a
vertex with only one remaining child is contracted.  Any finite
probability distribution can be represented in this way, by successively
splitting its support and using conditional probabilities.  The choice of
tree is part of the data; no associativity of matrix means is asserted.

\begin{lemma}[Properties of a finite mean tree]\label{transport:means}
Let $\mathcal T$ be a weighted binary tree with terminal weights $w_j$ and
positive definite inputs $A_j$, and let $M$ be its root.
The root commutes with invertible congruence, is monotone in each input,
and satisfies
\begin{equation}\label{transport:scalar-homogeneity}
 \mathcal T(c_jA_j:j)=\left(\prod_j c_j^{w_j}\right)M
 \qquad(c_j>0).
\end{equation}
For every unit vector $z$,
\begin{equation}\label{transport:vector-jensen}
 \langle z,Mz\rangle\leq\prod_j\langle z,A_jz\rangle^{w_j}.
\end{equation}
If $P$ is an orthogonal projection and $A_j\geq c_jP$ with $c_j>0$, then
\begin{equation}\label{transport:projection-transfer}
 M\geq\left(\prod_jc_j^{w_j}\right)P.
\end{equation}
Finally, suppose $A_j=\prod_{g=1}^K A_{j,g}$ with positive definite
factors, and every factor from band
$g$ commutes with every factor from band $g'$ whenever $g\ne g'$.
Then $M=\prod_{g=1}^K M_g$, where $M_g$ is the root of the same weighted
tree with inputs $A_{j,g}$.  The factors $M_g$ commute.
\end{lemma}

\begin{proof}
We first prove the assertions for one binary vertex.  For an invertible
$S$, put $\widehat A=S^*AS$ and
$U=A^{1/2}S\widehat A^{-1/2}$.  Then $U$ is unitary and
\[
 \widehat A^{-1/2}S^*BS\widehat A^{-1/2}
 =U^*(A^{-1/2}BA^{-1/2})U.
\]
Functional calculus proves
$(S^*AS)\#_p(S^*BS)=S^*(A\#_pB)S$.
After congruence to $A=I$, the identity
$A\#_pB=B\#_{1-p}A$ reduces to
$C^p=C^{1/2}(C^{-1})^{1-p}C^{1/2}$.
For $0<p<1$, the scalar formula
\begin{equation}\label{transport:power-integral}
 x^p=\frac{\sin(\pi p)}{\pi}
       \int_0^\infty \lambda^{p-1}x(x+\lambda)^{-1}\,d\lambda
\end{equation}
follows by substituting $\lambda=xt$ in the beta integral.
The inverse reverses the order of positive definite operators, so
$C(C+\lambda I)^{-1}=I-\lambda(C+\lambda I)^{-1}$ is increasing in $C$.
Thus $C\mapsto C^p$ is operator monotone.  This proves monotonicity in
$B$, and symmetry proves it in $A$.  The endpoint cases are immediate.
Scalar homogeneity at a vertex is
$(cA)\#_p(dB)=c^{1-p}d^p(A\#_pB)$.

For the vector inequality let $r=\langle z,Az\rangle$ and
$y=A^{1/2}z/\sqrt r$.  Scalar Jensen applied to the spectral measure of
$C=A^{-1/2}BA^{-1/2}$ in the unit vector $y$ gives
\[
 \langle z,(A\#_pB)z\rangle
 =r\langle y,C^py\rangle
 \leq r\langle y,Cy\rangle^p
 =\langle z,Az\rangle^{1-p}\langle z,Bz\rangle^p.
\]
Induction proves the congruence, monotonicity, homogeneity, and vector
inequality for the full tree.

To prove \eqref{transport:projection-transfer}, fix $0<\varepsilon<1$.
For every leaf,
\[
 A_j-(1-\varepsilon)c_jP
 =(1-\varepsilon)(A_j-c_jP)+\varepsilon A_j
 \geq\varepsilon\lambda_{\min}(A_j)I.
\]
Choose a common $\delta>0$ so small that
$(1-\varepsilon)c_j\delta\leq\varepsilon\lambda_{\min}(A_j)$ for all
leaves.  Monotonicity and homogeneity applied to the strictly positive
inputs $(1-\varepsilon)c_j(P+\delta I)$ give
\[
 M\geq(1-\varepsilon)\left(\prod_jc_j^{w_j}\right)(P+\delta I)
 \geq(1-\varepsilon)\left(\prod_jc_j^{w_j}\right)P.
\]
Let $\varepsilon$ decrease to zero.  This argument needs no definition of
the mean on singular inputs.

For the last assertion, cross-band commutation implies
\[
 \left(\prod_gD_g\right)\#_p\left(\prod_gE_g\right)
 =\prod_g(D_g\#_pE_g).
\]
Indeed the inverse square root of the first product factors, its
whitened second input is the product of the commuting positive operators
$D_g^{-1/2}E_gD_g^{-1/2}$, and the $p$th power factors.  Induction over the
tree completes the proof.
\end{proof}

We next describe how a perturbation at a leaf reaches the root.  A linear
map on matrices is completely positive if its entrywise action on every
matrix amplification preserves positive semidefinite matrices.  It is
unital if it maps the identity to the identity.

\begin{lemma}[Normalized derivative maps]\label{transport:cp}
At a binary vertex with output $M$ and a child $D$ of positive edge
weight $w$, define
\begin{equation}\label{transport:edge-map}
 \Phi_D(Z)=\frac1w M^{-1/2}
   \partial_D M[D^{1/2}ZD^{1/2}]M^{-1/2}.
\end{equation}
Here $\partial_D$ is the Fr\'echet derivative with the other child and the
weight fixed.  This map is unital and completely positive.  Composing
these maps from a terminal leaf $j$ to the root gives a unital completely
positive map $\Phi_j$ such that
\begin{equation}\label{transport:leaf-map}
 M^{-1/2}\partial_{A_j}M[A_j^{1/2}ZA_j^{1/2}]M^{-1/2}
 =w_j\Phi_j(Z).
\end{equation}
If all tree inputs are block diagonal for a fixed orthogonal
decomposition, these maps preserve each matrix block.  On the first
diagonal block they are the normalized derivative maps for the tree of
first diagonal blocks and do not depend on the other diagonal blocks.
\end{lemma}

\begin{proof}
Differentiating \eqref{transport:power-integral} gives
\begin{equation}\label{transport:power-derivative}
 D(C^p)[H]=\frac{\sin(\pi p)}{\pi}
 \int_0^\infty\lambda^p(C+\lambda I)^{-1}H(C+\lambda I)^{-1}\,d\lambda.
\end{equation}
The integral converges in operator norm: its integrand is $O(\lambda^p)$
at zero and $O(\lambda^{p-2})$ at infinity.  Each integrand is a
completely positive sandwich map.  Formula
\eqref{transport:power-derivative}, together with the congruences in
\eqref{transport:mean-def}, proves complete positivity for the second
child.  Exchange symmetry proves it for the first child.  Scalar
homogeneity gives $\partial_D M[D]=wM$, proving unitality.
The chain rule, with the square-root normalizations cancelling at
successive vertices, proves \eqref{transport:leaf-map}.  Composition
preserves both properties; the map along an empty path is the identity.
In the derivative integral all sandwich
factors are block diagonal when the inputs are block diagonal; their
first blocks depend only on the first blocks of the inputs.  The same
is true after exchanging the children and composing the maps.
\end{proof}

\subsection{Fourier formulas for inverse powers}

Write $\operatorname{ad}_Q(Z)=QZ-ZQ$.  The functions below convert
normalized matrix derivatives into derivatives of vector norms and into
energy errors.  Fix $0<s<1/2$ and set
\begin{equation}\label{transport:fourier-density}
 g_s(z)=\frac{\sinh(sz)}{\sinh(z/2)},\qquad
 m_s(u)=\frac{\sin(2\pi s)}{\cosh(2\pi u)+\cos(2\pi s)},
 \qquad g_s(0)=2s.
\end{equation}
The function $g_s$ is a rescaling of the classical positive-definite
hyperbolic-sine kernel studied in
\cite[Theorem~3.2]{BhatiaParthasarathy2000}.  We give the Fourier
normalization and the shifted-density comparison needed here.

\begin{lemma}[The two Fourier multipliers]\label{transport:fourier}
For real $z$,
\begin{equation}\label{transport:g-fourier}
 g_s(z)=\int_{\R}e^{iuz}m_s(u)\,du,
 \qquad m_s(u)>0,\qquad \int_{\R}m_s(u)\,du=2s.
\end{equation}
For $s=1/4$, the functions
\begin{equation}\label{transport:h-def}
 h_\pm(z)=\frac{e^{\pm sz}-1}{2\sinh(z/2)}
 =\pm e^{\pm sz/2}g_{s/2}(z)
\end{equation}
have Fourier densities
\begin{equation}\label{transport:q-density}
 q_\pm(u)=\pm m_{s/2}(u\pm is/2),\qquad
 |q_\pm(u)|\leq\frac1{\sqrt2}m_s(u).
\end{equation}
The quotients in \eqref{transport:h-def} are interpreted continuously
at zero.
\end{lemma}

\begin{proof}
For fixed real $z\ne0$, integrate $e^{iuz}m_s(u)$ around a rectangle
between the lines $\R$ and $\R+i$.  The vertical integrals tend to zero
because the denominator grows exponentially with $|\Re u|$.
The two poles in the strip are $i(1/2-s)$ and $i(1/2+s)$, with residues
of $m_s$ equal to $1/(2\pi i)$ and $-1/(2\pi i)$, respectively.
Since $m_s(u+i)=m_s(u)$, the residue theorem gives
\[
 (1-e^{-z})\int_{\R}e^{iuz}m_s(u)\,du
 =e^{-(1/2-s)z}-e^{-(1/2+s)z}.
\]
Division gives \eqref{transport:g-fourier}; dominated convergence gives
its value at zero.  Positivity follows from the displayed denominator.

The algebraic identities in \eqref{transport:h-def} follow by factoring
$e^{\pm sz}-1$.  The nearest poles of $m_{s/2}$ have imaginary parts
of absolute value $1/2-s/2$, larger than $s/2$.  Shifting its Fourier
contour by $\pm is/2$ therefore crosses no pole and gives the densities
in \eqref{transport:q-density}.  To check the bound without losing a
constant, write $c=\cosh(2\pi u)$.  At $s=1/4$,
\[
 m_s(u)=c^{-1},\qquad
 m_{s/2}(u\pm is/2)=\frac1{c+1\pm i\sinh(2\pi u)}.
\]
Consequently
\[
 \frac{|m_{s/2}(u\pm is/2)|}{m_s(u)}
 =\sqrt{\frac{c}{2(c+1)}}\leq\frac1{\sqrt2}.
\]
\end{proof}

\subsection{Finite histories and entropy gains}

We specify the data to which the transport estimate applies.  Let
$\mathcal V$ be a fixed finite tensor product of one-copy physical and
auxiliary spaces, and put $\mathcal S_k=\Sym^k\mathcal V$.  The metrics
$W_{Q,k}$ are defined in \eqref{replicas:w-definition} and studied in
Lemma~\ref{replicas:metric}, with the common padding dimension
$\dim\mathcal V$.  We suppress $k$ in $W_{Q,k}$.  A partition
$(P,Y,F)$ of these tensor factors has band metric
\begin{equation}\label{transport:band-metric}
 A(P,Y,F)=(W_P^{-1}W_F^{-1}W_Y)^2
 \quad\hbox{on }\mathcal S_k.
\end{equation}
All three factors in this formula commute and are positive definite.

Fix a positive integer $K$ and a finite weighted history tree with leaves
$h\in\mathcal H$ and weights $w_h>0$.  For each $h$ and
$g\in\{1,\ldots,K\}$ choose an old partition
$(P_{h,g},Y_{h,g},F_{h,g})$.  At each history $h$ also fix a finite
conditional choice tree with leaves $c\in\mathcal C_h$ and weights
$q_{c\mid h}>0$.  A choice specifies a new partition in every band.
Each new partition is obtained either by doing nothing or by transferring
a tensor subsystem $x\subseteq Y$ entirely to $P$ or entirely to $F$.
Write $A_{h,g}$ and $A_{h,c,g}$ for the old and new band metrics and set
\[
 A_h=\prod_{g=1}^K A_{h,g},\qquad
 A_{h,c}=\prod_{g=1}^K A_{h,c,g}.
\]
We impose the following commutation condition: every old or new metric
from band $g$ commutes with every old or new metric from band $g'$ for
$g\ne g'$, including metrics at different histories and choices.
No commutation within a band is required.

Let $B_h$ be the conditional choice-tree root with inputs $A_{h,c}$.
For one common parameter $p\in[0,1]$, replace each old input of the
history tree by $A_h\#_pB_h$ and call the resulting root
$\mathcal M(p)$.  Its terminal leaves at $0<p<1$ are
\begin{equation}\label{transport:terminal-weights}
 j=(h,\mathrm{old}),\quad \pi_j=(1-p)w_h,
 \qquad\hbox{or}\qquad
 j=(h,c,\mathrm{new}),\quad \pi_j=pw_hq_{c\mid h}.
\end{equation}
Their metrics are $A_h$ and $A_{h,c}$, respectively.  At $p=1$ this is
the history tree with the conditional choice trees attached; it can serve
as the old tree of a subsequent round.  The same $p$ is used at every
history in a round.

Figure~\ref{transport:tree-figure} displays the two types of terminal
leaf.  The history and conditional-choice weights are products along
paths; the single interpolation vertex contributes $1-p$ or $p$.
\begin{figure}[H]
\centering
\begin{tikzpicture}[x=1cm,y=1cm,font=\small,>=Stealth]
 \node (root) at (0,0) {$\mathcal M(p)$};
 \node[align=center] (node) at (4,0) {$A_h\#_p B_h$};
 \node (old) at (7,1.2) {$A_h$};
 \node (choice) at (7,-1.2) {$B_h$};
 \node (new) at (11,-1.2) {$A_{h,c}$};
 \draw[->,densely dotted] (node)--node[above,align=center]
   {history path\\product $w_h$}(root);
 \draw[->] (old)--node[above left] {$1-p$}(node);
 \draw[->] (choice)--node[below left] {$p$}(node);
 \draw[->,densely dotted] (new)--node[above,align=center]
   {choice path\\product $q_{c\mid h}$}(choice);
 \node[above=3pt of old] {$\pi_{(h,\mathrm{old})}=(1-p)w_h$};
 \node[below=3pt of new] {$\pi_{(h,c,\mathrm{new})}=pw_hq_{c\mid h}$};
\end{tikzpicture}
\caption{One history and one conditional choice in a round.  Arrows
point toward the root and show metric evaluation; every internal
vertex uses its specified geometric mean.  Dotted arrows abbreviate
whole paths, not single edges.  The same interpolation parameter $p$
is used at every history.  Trace adjoints of normalized derivatives
carry states in the reverse direction.}
\label{transport:tree-figure}
\end{figure}

For a unit vector $\theta\in\mathcal V$, the entropy assigned to a move
of $x$ to $P$, with $Y=xY_0$, is
\begin{equation}\label{transport:move-eta}
 \eta_{h,c,g}(\theta)
 =S_\theta(x\mid P)+S_\theta(x\mid Y_0)
 =I_\theta(x:F\mid P)\geq0.
\end{equation}
For a move to $F$ exchange $P$ and $F$ in this definition.  A choice
that makes no move has $\eta_{h,c,g}=0$.  The equality in
\eqref{transport:move-eta} follows by purity on $PxY_0F$.
Choose $\ell\geq1$ bounding $\log(e\dim x)$ for every transferred
subsystem.  Throughout, $a=2t>0$ satisfies the smallness hypotheses of
Lemmas~\ref{lem:relative-pin} and \ref{lem:skew}; in particular
$t<1/4$ and $a\ell$ is bounded by a sufficiently small absolute constant.

For the energy part of the estimate, let $0\leq h_i\leq I$, $i$ in a
finite index set, act on physical tensor factors, and let $D_i$ be a
designated tensor support containing the actual support of $h_i$.
For each terminal partition and each band, assume $D_i$ is contained in
one of its three parts, except possibly in one band.  In that exceptional
band, $D_i$ meets $Y$ and exactly one of $P,F$.  Call this a split of
$i$ at leaf $j$, and write $\mathcal J_i$ for the set of such leaves.
For a split put $x=D_i\cap Y$, let the receiving side be the part
$P$ or $F$ met by $D_i$, and define $\eta_{i,j}$ by the corresponding
formula \eqref{transport:move-eta}.  In addition to the bounds for all
transferred subsystems, require
$\log(e\dim D_i)\leq\ell$ for every $i$ with
$\mathcal J_i\ne\varnothing$, retaining the smallness condition on
$a\ell$.  No support-dimension bound is required for terms that are
unsplit at every terminal leaf.
These are explicit assumptions on a finite collection of partitions and
supports; geometric ways to satisfy them are considered later.

Define
\begin{equation}\label{transport:replica-energy}
 \bar h_i=\left.\frac1k\sum_{r=1}^k h_{i,r}\right|_{\mathcal S_k},
 \qquad \bar H=\sum_i\bar h_i,
 \qquad W_i(p)=\sum_{j\in\mathcal J_i}\pi_j.
\end{equation}
If a leaf does not split $i$, its metric commutes with $\bar h_i$.
Indeed, when $D_i$ is contained in a partition part, the sum of its
copy-translates commutes with every copy permutation on that part and
with every permutation on its disjoint complement.  It therefore
commutes with the Schur-label functions defining that band's metric.
Apply this argument in every band.

For use in the statement, recall the coherent vectors and measures
\begin{equation}\label{transport:coherent-measure}
 P_{\theta,k}=|\theta^{\otimes k}\rangle\langle\theta^{\otimes k}|,
 \quad D_k=\dim\mathcal S_k,\quad
 d\mu_\sigma(\theta)=D_k\Tr(\sigma P_{\theta,k})\,d\theta,
\end{equation}
where $d\theta$ is Haar probability on the unit sphere of $\mathcal V$.
For a density matrix $\sigma$ on $\mathcal S_k$, this is a probability
measure, since $D_k\int P_{\theta,k}\,d\theta=I_{\mathcal S_k}$.

\begin{proposition}[Entropy and energy transport]\label{prop:transport}
Fix a finite one-copy space, history and conditional choice trees,
partitions, cross-band commutation, positive parameter $a$, and positive
contractions with designated supports as specified above.  Let $\ell$
bound the logarithmic dimensions of every transferred subsystem and
every designated support $D_i$ for which $\mathcal J_i\ne\varnothing$,
with the stipulated smallness of $a\ell$.
For every $k$ and nonzero $\mathrm{pre}\in\mathcal S_k$, set
\begin{equation}\label{transport:filtered-vector}
 s=\frac14,\qquad
 N(p)=\norm{\mathcal M(p)^{-s}\mathrm{pre}},\qquad
 v(p)=\frac{\mathcal M(p)^{-s}\mathrm{pre}}{N(p)}.
\end{equation}
For every $0<p<1$, let $\Phi_j$ be the normalized map from terminal
leaf $j$ to $\mathcal M(p)$ and define density matrices
\begin{equation}\label{transport:states}
 \sigma_{j,u}=\Phi_j^*\bigl(
 |\mathcal M(p)^{-iu}v(p)\rangle
 \langle\mathcal M(p)^{-iu}v(p)|\bigr),\qquad u\in\R.
\end{equation}
The adjoint here is for the trace pairing:
$\Tr(\Phi_j^*(\sigma)Z)=\Tr(\sigma\Phi_j(Z))$.
Write $\sigma_{h,u}$ for the state at $j=(h,\mathrm{old})$, and put
$C_h=A_h^{-1/2}B_hA_h^{-1/2}$.  Then
\begin{equation}\label{transport:exact-derivative}
 -\partial_p\log N(p)^2
 =\sum_h w_h\int_{\R}m_s(u)\Tr(\sigma_{h,u}\log C_h)\,du.
\end{equation}
There is $\beta_k=O_{\mathrm{fixed}}(\log(k+1))$, independent of $p$
and $\mathrm{pre}$, such that
\begin{align}\label{transport:entropy-gain}
 -\partial_p\log N(p)^2
 \geq{}&ka\sum_{h,g}w_h\int_{\R}m_s(u)
    \int\sum_{c\in\mathcal C_h}q_{c\mid h}
          \eta_{h,c,g}(\theta)\,d\mu_{\sigma_{h,u}}(\theta)\,du
 \notag\\[-2pt]
 &{}-CkaK a^{1/4}\ell^C-\beta_k.
\end{align}
If in addition $\bar H\,\mathrm{pre}=E_0\mathrm{pre}$, then
\begin{equation}\label{transport:energy}
 \langle v,\bar Hv\rangle
 \leq 2E_0+C a^2\ell^C
 \sum_i W_i(p)\sum_{j\in\mathcal J_i}\pi_j
 \int_{\R}m_s(u)\int\eta_{i,j}(\theta)^{1/8}
                 \,d\mu_{\sigma_{j,u}}(\theta)\,du+o_k(1).
\end{equation}
The remainder in \eqref{transport:energy} is uniform in $p\in(0,1)$,
in nonzero $\mathrm{pre}$ satisfying the eigenvector assumption, and in
all replica states.  The constants $C$ and the exponents denoted by $C$
in \eqref{transport:entropy-gain} and \eqref{transport:energy} are
universal; support dimensions enter those coefficients only through
$\ell$.  Only the implicit constant in
$\beta_k=O_{\mathrm{fixed}}(\log(k+1))$ and the convergence rate of
$o_k(1)$ may depend on the entire fixed finite data.  Neither depends
on $p$ or the replica state.  The entropy estimate can be integrated
over a closed subinterval
of $[0,1]$ using the continuous endpoint values of $N$.
\end{proposition}

\begin{proof}
We first identify the derivative states and prove the entropy estimate.
We then construct an auxiliary block tree for each physical term and
derive \eqref{transport:energy} by a matrix Schwarz inequality.

\paragraph{The derivative at an interpolating vertex.}
At the vertex $M_h=A_h\#_pB_h$, put
$V_h=C_h^{p/2}A_h^{1/2}M_h^{-1/2}$.  Since $V_h^*V_h=I$, this matrix
is unitary, and direct differentiation gives
\begin{equation}\label{transport:node-p-derivative}
 M_h^{-1/2}\partial_pM_hM_h^{-1/2}=V_h^*(\log C_h)V_h.
\end{equation}
This is also the normalized old-child derivative applied to $\log C_h$.
To verify that assertion, perturb the old input to
$A_h^{1/2}(I+\delta\log C_h)A_h^{1/2}$ while keeping $B_h$ fixed.
After congruence the two inputs commute, so their mean is
\[
 (I+\delta\log C_h)\#_pC_h
 =(I+\delta\log C_h)^{1-p}C_h^p.
\]
Its derivative at zero is $(1-p)C_h^p\log C_h$.  The normalization in
\eqref{transport:edge-map} divides by $1-p$ and proves the assertion.
Composing with the ancestor maps therefore yields
\begin{equation}\label{transport:root-p-derivative}
 \mathsf H:=\mathcal M^{-1/2}(\partial_p\mathcal M)\mathcal M^{-1/2}
 =\sum_h w_h\Phi_{(h,\mathrm{old})}(\log C_h).
\end{equation}
The coefficient here is $w_h$, whereas the weight of the same terminal
leaf in the full tree is $(1-p)w_h$.

To differentiate the norm, diagonalize $\mathcal M$ with eigenvalues
$x_r>0$ and substitute $\mathrm{pre}=N\mathcal M^s v$.
The matrix coefficient, between the components $\overline v_r$ and
$v_t$, multiplying the $(r,t)$ entry of $\mathsf H$ is
\[
 -(x_rx_t)^{s+1/2}
       \frac{x_r^{-2s}-x_t^{-2s}}{x_r-x_t}
 =g_s(\log x_r-\log x_t),
\]
with the divided difference interpreted as a derivative when the
eigenvalues agree.  The common value in that case is $2s$.
Consequently
\begin{equation}\label{transport:norm-derivative}
 -\partial_p\log N^2
 =\langle v,g_s(\operatorname{ad}_{\log\mathcal M})(\mathsf H)v\rangle.
\end{equation}
Insert \eqref{transport:root-p-derivative} and the Fourier formula
\eqref{transport:g-fourier}.  Lemma~\ref{transport:cp} implies that
\eqref{transport:states} is positive and has trace one.  Moving each
map to its trace adjoint proves \eqref{transport:exact-derivative}.

\paragraph{Transfer of the relative metric lower bound.}
By Lemma~\ref{transport:means}, the old metrics, the conditional roots,
the interpolated root $\mathcal M=\prod_g\mathcal M_g$, and the
relative metrics factor over bands.  Here $\mathcal M_g$ is the full
tree evaluated on the inputs from band $g$ alone.  More precisely, if
$B_{h,g}$ is the conditional tree of new metrics in band $g$, then
\begin{equation}\label{transport:relative-factorization}
 C_h=\prod_g C_{h,g},\qquad
 C_{h,g}=A_{h,g}^{-1/2}B_{h,g}A_{h,g}^{-1/2},\qquad
 \log C_h=\sum_g\log C_{h,g}.
\end{equation}
Congruence covariance identifies $C_{h,g}$ as the conditional choice
tree of the single-move relative metrics
$A_{h,g}^{-1/2}A_{h,c,g}A_{h,g}^{-1/2}$.
Lemma~\ref{lem:relative-pin} gives, for every unit $\theta$,
\[
 A_{h,g}^{-1/2}A_{h,c,g}A_{h,g}^{-1/2}
 \geq b_k\exp\!\left(ka\left[
       \eta_{h,c,g}(\theta)-C a^{1/4}\ell^C\right]\right)P_{\theta,k},
\]
where $b_k>0$ and $-\log b_k=O_{\mathrm{fixed}}(\log(k+1))$.
The same $b_k$ and constant can be used for the finitely many moves.
For a move of the empty subsystem the relative metric is the identity,
and the assertion holds after decreasing $b_k$ to at most one.
Projection transfer, \eqref{transport:projection-transfer}, gives
\begin{equation}\label{transport:conditional-pin}
 C_{h,g}\geq b_k
 \exp\!\left(ka\left[\sum_cq_{c\mid h}\eta_{h,c,g}(\theta)
                    -C a^{1/4}\ell^C\right]\right)P_{\theta,k}.
\end{equation}
The common polynomial factor occurs once, because the conditional
weights sum to one.

Here is the logarithmic passage from this pointwise rank-one inequality
to an estimate valid in every state.  Define the unital positive map
\[
 \mathcal Q_k(f)=D_k\int f(\theta)P_{\theta,k}\,d\theta.
\]
The identity $\mathcal Q_k(1)=I$ follows from the invariance of the
integral and irreducibility of $\Sym^k\mathcal V$ under the full unitary
group; its trace fixes the scalar.  For a positive bounded function
$f$ bounded away from zero, integration of the positive scalar block
matrix with entries $f,1;1,f^{-1}$ gives
\[
 \begin{pmatrix}\mathcal Q_k(f)&I\\I&\mathcal Q_k(f^{-1})\end{pmatrix}
 \geq0,\qquad
 \mathcal Q_k(f)^{-1}\leq\mathcal Q_k(f^{-1}).
\]
Apply this also to $f+\lambda$ and integrate the scalar resolvent
formula
\[
 \log x=\int_0^\infty
       \left((1+\lambda)^{-1}-(x+\lambda)^{-1}\right)d\lambda.
\]
It proves
\begin{equation}\label{transport:log-jensen}
 \log\mathcal Q_k(f)\geq\mathcal Q_k(\log f).
\end{equation}
The same inverse-order argument proves operator monotonicity of the
logarithm.  All these integrals converge in norm on compact positive
spectral intervals.

Integrating \eqref{transport:conditional-pin} against Haar probability
gives
\[
 C_{h,g}\ge\frac{b_k}{D_k}\,
 \mathcal Q_k\!\left(\exp\!\left\{ka\left[
    \sum_cq_{c\mid h}\eta_{h,c,g}
                    -C a^{1/4}\ell^C\right]\right\}\right).
\]
Thus operator log monotonicity and \eqref{transport:log-jensen} turn
the factor $D_k$ into an additive logarithmic loss:
\begin{align}\label{transport:log-relative}
 \log C_{h,g}\geq{}&ka\mathcal Q_k\!\left(
       \sum_cq_{c\mid h}\eta_{h,c,g}-C a^{1/4}\ell^C\right)
       +(\log b_k-\log D_k)I.
\end{align}
The entropies are bounded in the fixed one-copy dimension, so the
exponential symbol used here is positive and bounded away from zero
for each fixed $k$.  Also
$D_k=\binom{k+\dim\mathcal V-1}{\dim\mathcal V-1}$ is polynomial in
$k$.  Taking the expectation in the trace-one state $\sigma_{h,u}$,
summing \eqref{transport:relative-factorization}, and using
$\sum_hw_h=1$ and $\int m_s=1/2$ proves
\eqref{transport:entropy-gain}.  In particular the loss is logarithmic
in $k$, without an additional factor $D_k$ after taking expectations.

We have obtained the entropy gain against the old-leaf states.  To
control the energy in these same states, the remaining argument
represents the commutator with each $\bar h_i$ by a block matrix and
transports its square through the same tree.

\paragraph{An auxiliary positive metric for each energy term.}
Fix $i$.  Let $T:\mathcal S_k\longrightarrow\mathcal D$ be
$h_i^{1/2}$ on the last copy, restricted to $\mathcal S_k$, with
codomain its image.  Then
$T^*T=\bar h_i$: matrix elements of any marked copy coincide between
symmetric vectors, and their average is the restriction in
\eqref{transport:replica-energy}.  If $T=0$, the term contributes
nothing and can be omitted.  Otherwise $T$ is surjective onto
$\mathcal D$.

We will add a second diagonal block to each leaf metric.  This block
is chosen so that the squared off-diagonal commutator recovers the
absolute-value expression of Lemma~\ref{lem:symbol}.  The first block
remains the original metric, so Lemma~\ref{transport:cp} will preserve
its derivative maps when the enlarged tree is evaluated.
For every terminal leaf with input $A_j$, we construct a positive
definite operator $B_{i,j}$ on $\mathcal D$.  When
$[A_j,\bar h_i]=0$, use the polar decomposition $T=U|T|$.
Here $U$ is unitary from $\supp|T|$ onto $\mathcal D$, and this support
is invariant under $A_j$.  Set
$B_{i,j}=U(A_j|_{\supp|T|})U^*$.  Commutation with $|T|$ gives
\begin{equation}\label{transport:unsplit-block}
 Y_{i,j}:=B_{i,j}^{1/2}TA_j^{-1/2}=T
 =B_{i,j}^{-1/2}TA_j^{1/2}=:Z_{i,j}.
\end{equation}

For a split leaf put
\[
 O_{i,j}=A_j^{-1/2}\bar h_iA_j^{1/2},\qquad
 S_j=TA_j^{-1/2},\qquad Q_j=|O_{i,j}^*|.
\]
For any operator $O$, the notation $|O|$ means $(O^*O)^{1/2}$.
The range of $O_{i,j}$ is the range of $S_j^*$, namely
$A_j^{-1/2}\operatorname{ran}T^*$.  Since $S_j$ is surjective,
$S_jS_j^*$ is positive definite on $\mathcal D$.  Define
\begin{equation}\label{transport:rank-deficient-B}
 B_{i,j}=(S_jS_j^*)^{-1}S_jQ_jS_j^*(S_jS_j^*)^{-1}.
\end{equation}
The operator $Q_j$ is strictly positive on $\operatorname{ran}S_j^*$.
Thus \eqref{transport:rank-deficient-B} is positive definite on
$\mathcal D$ and $S_j^*B_{i,j}S_j=Q_j$; the latter follows because
$S_j^*(S_jS_j^*)^{-1}S_j$ is the projection onto
$\operatorname{ran}S_j^*$.  Define $Y_{i,j},Z_{i,j}$ by the formulas
in \eqref{transport:unsplit-block}.  They now satisfy
\[
 Y_{i,j}^*Y_{i,j}=|O_{i,j}^*|,
 \qquad Y_{i,j}^*Z_{i,j}=O_{i,j}.
\]
The map $Y_{i,j}$ is surjective.  Consequently
$Y_{i,j}(Y_{i,j}^*Y_{i,j})^+Y_{i,j}^*=I_{\mathcal D}$, where the plus
denotes the inverse on the support and zero on the kernel.  Multiplying
on the two sides by $Z_{i,j}^*$ and $Z_{i,j}$ shows
\[
 Z_{i,j}^*Z_{i,j}
 =O_{i,j}^*|O_{i,j}^*|^+O_{i,j}=|O_{i,j}|.
\]
The last equality follows, for example, from a singular value
decomposition of $O_{i,j}$.  Hence
\begin{equation}\label{transport:skew-square}
 \mathsf D_{i,j}:=(Y_{i,j}-Z_{i,j})^*(Y_{i,j}-Z_{i,j})
 =|O_{i,j}^*|+|O_{i,j}|-O_{i,j}-O_{i,j}^*\geq0.
\end{equation}
For an unsplit leaf use $\mathsf D_{i,j}=0$; it agrees with the same
formula because $O_{i,j}=\bar h_i$.  The construction has therefore
included singular $\bar h_i$ without ever inverting it on its kernel.

\paragraph{Transporting the square with both Fourier kernels.}
Evaluate the full terminal tree on the block inputs
\[
 M_{i,j}=\begin{pmatrix}A_j&0\\0&B_{i,j}\end{pmatrix}.
\]
Its root is
$\mathbb M_i=\operatorname{diag}(\mathcal M,\mathcal B_i)$.
Denote its normalized leaf maps by $\widehat\Phi_{i,j}$.  Their first
diagonal-block restrictions are exactly $\Phi_j$, by
Lemma~\ref{transport:cp}, although the second blocks depend on $i$.
Let
\[
 L_i=\begin{pmatrix}0&T^*\\T&0\end{pmatrix},\qquad
 K_{i,j}=M_{i,j}^{1/2}L_iM_{i,j}^{-1/2}
              -M_{i,j}^{-1/2}L_iM_{i,j}^{1/2}.
\]
Rotating every input by $e^{itL_i}$ rotates the root by the same unitary,
by congruence covariance.  Differentiate at zero, normalize at the
root, and cancel the common factor $-i$.  The chain rule gives
\begin{equation}\label{transport:block-rotation}
 2\sinh(\operatorname{ad}_{\log\mathbb M_i}/2)L_i
 =\sum_j\pi_j\widehat\Phi_{i,j}(K_{i,j}).
\end{equation}
The lower-left block of $K_{i,j}$ is $Y_{i,j}-Z_{i,j}$, its upper-right
block is the negative adjoint, and the upper-left block of
$K_{i,j}^*K_{i,j}$ is $\mathsf D_{i,j}$.  In particular
$K_{i,j}=0$ for $j\notin\mathcal J_i$.

Apply $h_\pm(\operatorname{ad}_{\log\mathbb M_i})$ to
\eqref{transport:block-rotation}.  By \eqref{transport:h-def}, the
lower-left block of the left side is
\begin{equation}\label{transport:E-def}
 E_{i,\pm}=\mathcal B_i^{\pm s}T\mathcal M^{\mp s}-T.
\end{equation}
The Fourier formula for the right side and preservation of blocks give
an expression for $(0,E_{i,\pm}v)$ as a sum over split leaves of
\[
 \pi_j\int_{\R}q_\pm(u)\,
 \mathbb M_i^{iu}\widehat\Phi_{i,j}(K_{i,j})
 \mathbb M_i^{-iu}(v,0)\,du.
\]
The output phase is unitary.  Cauchy--Schwarz with respect to the
positive measure assigning mass $\pi_j|q_\pm(u)|\,du$ to $(j,u)$
therefore bounds the squared norm by
\begin{align*}
 &W_i\left(\int_{\R}|q_\pm(u)|\,du\right)
 \sum_{j\in\mathcal J_i}\pi_j\int_{\R}|q_\pm(u)|
 \norm{\widehat\Phi_{i,j}(K_{i,j})(\mathcal M^{-iu}v,0)}^2\,du.
\end{align*}
For a unital completely positive map $\Psi$,
\begin{equation}\label{transport:cp-schwarz}
 \Psi(K)^*\Psi(K)\leq\Psi(K^*K).
\end{equation}
Indeed apply its two-by-two amplification to the positive block matrix
$\left(\begin{smallmatrix}K^*K&K^*\\K&I\end{smallmatrix}\right)$ and
take the Schur complement of the identity.  By the first-block
restriction of $\widehat\Phi_{i,j}$, this bounds the norm squared in
the preceding display by $\Tr(\sigma_{j,u}\mathsf D_{i,j})$.
Finally \eqref{transport:q-density} and $\int m_s=1/2$ give the exact
estimate
\begin{equation}\label{transport:block-bound}
 \norm{E_{i,\pm}v}^2
 \leq\frac14 W_i\sum_{j\in\mathcal J_i}\pi_j
       \int_{\R}m_s(u)\Tr(\sigma_{j,u}\mathsf D_{i,j})\,du.
\end{equation}
This proves the bound for both signs.  The first factor $W_i$ comes
from Cauchy--Schwarz, and the sum contains a second total split weight.
For example, a uniform bound on its trace integrals gives a bound
proportional to $W_i^2$, rather than $W_i$.

\paragraph{From the transported square to physical energy.}
At a split leaf only one band fails to commute with $\bar h_i$.
The other factors cancel in its similarity transform $O_{i,j}$.
At such a leaf $j\in\mathcal J_i$, the assumed bound
$\log(e\dim D_i)\leq\ell$ applies.  Our designated-support assumptions
put this two-part split within the hypotheses of
Lemma~\ref{lem:symbol}; Lemma~\ref{lem:skew}
then gives, uniformly over density matrices $\sigma$ on $\mathcal S_k$,
\begin{equation}\label{transport:symbol-cost}
 \Tr(\sigma\mathsf D_{i,j})
 \leq C a^2\ell^C\int\eta_{i,j}(\theta)^{1/8}
               \,d\mu_\sigma(\theta)+r_k,
 \qquad r_k\longrightarrow0.
\end{equation}
There are finitely many split term--leaf pairs, so one common
$r_k\geq0$ can be chosen for all of them; take $r_k=0$ if there are none.
Unsplit pairs have $\mathsf D_{i,j}=0$ exactly and incur no symbol
remainder or support-dimension factor.  These terminal metrics do not depend
on the interpolation parameter.  Uniformity over $\sigma$ makes
\eqref{transport:symbol-cost} applicable to every $\sigma_{j,u}$,
regardless of $p$, $u$, or the normalized vector.

The two opposite powers of $\mathcal B_i$ in
\eqref{transport:E-def} cancel in the pairing:
\begin{align}\label{transport:energy-similarity}
 \Re\langle(T+E_{i,+})v,(T+E_{i,-})v\rangle
 &=\Re\langle\mathcal M^{-s}v,\bar h_i\mathcal M^s v\rangle
 \notag\\
 &\geq\frac12\norm{Tv}^2
       -\frac32\bigl(\norm{E_{i,+}v}^2+\norm{E_{i,-}v}^2\bigr).
\end{align}
For the last inequality, expand the left pairing, bound each linear
cross term by $\frac14\norm{Tv}^2+\norm{E_{i,\pm}v}^2$, and bound
the remaining mixed term by half the sum of the two error squares.
Summing the equality in \eqref{transport:energy-similarity} over $i$
gives $E_0$: one has $\mathcal M^s v=\mathrm{pre}/N$, and
$\langle\mathcal M^{-s}v,\mathcal M^s v\rangle=1$.
Thus
\[
 \langle v,\bar Hv\rangle
 \leq 2E_0+3\sum_{i,\pm}\norm{E_{i,\pm}v}^2.
\]
Insert \eqref{transport:block-bound} and
\eqref{transport:symbol-cost} to obtain \eqref{transport:energy}.
The total remainder is bounded by a fixed constant times
$r_k\sum_iW_i^2$ and hence tends to zero uniformly in $p$ and $v$.

Finally, all tree inputs are positive definite, so for fixed $k$ the
roots and $N$ are smooth up to $p=0,1$.  The normalized maps are used
only where their branch weights are positive; their being unital
ensures bounds uniform as $p$ approaches an endpoint.  The entropies
are bounded in the fixed one-copy space, and all probability weights
sum to one.  Thus the entropy integrand and its error have the uniform
bounds required to integrate \eqref{transport:entropy-gain} and pass to
the endpoints.  A finite sequence of rounds permits the same choices
of polynomial losses and remainders by taking their maxima.
\end{proof}

The proposition separates two uses of a history.  Its old-leaf state
in \eqref{transport:exact-derivative} is exactly the state used for that
leaf in \eqref{transport:block-bound}.  Only the scalar weights differ:
$w_h$ in the derivative and $(1-p)w_h$ in the energy estimate.  This
identity allows a later averaging argument to select one common time
and one common interpolation parameter while retaining control of all
energy terms.
